% Created 2026-04-24 sex 22:59
% Intended LaTeX compiler: xelatex
\documentclass[11pt]{article}
\usepackage{graphicx}
\usepackage{longtable}
\usepackage{wrapfig}
\usepackage{rotating}
\usepackage[normalem]{ulem}
\usepackage{capt-of}
\usepackage{hyperref}
\date{2026}
\title{Lab 05 — Predição de Preços de Imóveis\\\medskip
\large Documentação da Solução Vencedora}
\hypersetup{
 pdfauthor={},
 pdftitle={Lab 05 — Predição de Preços de Imóveis},
 pdfkeywords={},
 pdfsubject={},
 pdfcreator={Emacs 30.2 (Org mode 9.7.11)}, 
 pdflang={English}}
\begin{document}

\maketitle
\setcounter{tocdepth}{2}
\tableofcontents

\section{Visão Geral}
\label{sec:org18ea467}

O problema consiste em prever o preço de venda de imóveis residenciais
(dataset Ames Housing) a partir de 79 variáveis descritivas. A métrica
de avaliação é o RMSE calculado sobre \texttt{log(SalePrice)}.

\begin{center}
\begin{tabular}{lll}
Métrica & Baseline (regressão linear) & Solução final\\
\hline
RMSE\_LOG OOF & \textasciitilde{}0.23 & 0.1140\\
Posição & — & 1º lugar\\
\end{tabular}
\end{center}

A solução é um \textbf{stacking ensemble} de dois níveis: cinco modelos de
gradient boosting como camada base e um blend de RidgeCV + ElasticNet
como meta-modelo.
\section{Pré-processamento}
\label{sec:org332e0c5}

\subsection{Remoção de Outliers}
\label{sec:orgfdda29f}

Antes de qualquer modelagem, quatro filtros são aplicados ao conjunto
de treino para eliminar observações atípicas que distorceriam o
aprendizado:

\begin{verbatim}
train_df = train_df[train_df["GrLivArea"] < 4500]
train_df = train_df[~((train_df["GrLivArea"] > 4000) & (train_df["SalePrice"] < 200000))]
train_df = train_df[train_df["LotArea"] < 100000]
train_df = train_df[train_df["TotalBsmtSF"] < 6000]
\end{verbatim}

O segundo filtro é especialmente importante: remove casas grandes
vendidas a preço baixo, que são anomalias conhecidas do dataset Ames e
prejudicam modelos de árvore.
\subsection{Transformação da Variável Alvo}
\label{sec:org025655d}

O target \texttt{SalePrice} é transformado com \texttt{log1p} antes do treino:

\begin{verbatim}
y = np.log1p(train_df["SalePrice"])
\end{verbatim}

Isso aproxima a distribuição da normal, reduz a influência de valores
extremos e é o formato esperado na submissão — as predições são
entregues em escala logarítmica, sem reverter com \texttt{expm1}.
\subsection{Imputação Contextual}
\label{sec:org88a79f4}

Em vez de imputação genérica por mediana, os valores ausentes são
tratados com base no domínio do problema:

\begin{itemize}
\item \textbf{Categóricas de ausência semântica} (ex: \texttt{GarageType}, \texttt{PoolQC}): NaN
significa que o imóvel não possui aquela característica, portanto
preenchidos com a string \texttt{"None"}.

\item \textbf{Numéricas de área ou contagem} (ex: \texttt{GarageArea}, \texttt{BsmtFinSF1}):
imóveis sem garagem ou porão têm área zero, portanto preenchidos
com \texttt{0}.

\item \textbf{Demais numéricas}: imputação por mediana dentro de cada fold de
validação cruzada, usando \texttt{SimpleImputer}.
\end{itemize}
\section{Feature Engineering}
\label{sec:org287a936}

A função \texttt{add\_features} expande o conjunto original de 80 para 103
features, agrupadas em quatro categorias:
\subsection{Agregações de Área}
\label{sec:org4e94e8e}

\begin{verbatim}
df["TotalSF"]    = df["TotalBsmtSF"] + df["1stFlrSF"] + df["2ndFlrSF"]
df["TotalPorchSF"] = df["OpenPorchSF"] + df["EnclosedPorch"] + ...
df["TotalBathrooms"] = df["FullBath"] + 0.5*df["HalfBath"] + ...
\end{verbatim}

Capturam a área utilizável total, que é mais preditiva que qualquer
área individual.
\subsection{Flags Binárias de Presença}
\label{sec:org85b3ca8}

\begin{verbatim}
df["HasPool"]      = (df["PoolArea"] > 0).astype(int)
df["HasGarage"]    = (df["GarageArea"] > 0).astype(int)
df["HasBsmt"]      = (df["TotalBsmtSF"] > 0).astype(int)
df["HasFireplace"] = (df["Fireplaces"] > 0).astype(int)
df["IsRemodeled"]  = (df["YearRemodAdd"] != df["YearBuilt"]).astype(int)
df["IsNew"]        = (df["YrSold"] == df["YearBuilt"]).astype(int)
\end{verbatim}

Modelos de árvore beneficiam-se de splits binários explícitos quando a
distinção presença/ausência é mais relevante que a magnitude.
\subsection{Interações Multiplicativas}
\label{sec:orgb17fea7}

\begin{verbatim}
df["OverallQual_TotalSF"]   = df["OverallQual"] * df["TotalSF"]
df["OverallQual_GrLivArea"] = df["OverallQual"] * df["GrLivArea"]
df["OverallQual_sq"]        = df["OverallQual"] ** 2
df["QualCondScore"]         = df["OverallQual"] * df["OverallCond"]
\end{verbatim}

\texttt{OverallQual} é a variável mais correlacionada com preço. Suas
interações com área capturam o efeito não-linear de que qualidade
importa proporcionalmente mais em casas grandes.
\subsection{Transformações Logarítmicas e Temporais}
\label{sec:orgc834dca}

\begin{verbatim}
df["LotAreaLog"]  = np.log1p(df["LotArea"])
df["GrLivAreaLog"]= np.log1p(df["GrLivArea"])
df["HouseAge"]    = df["YrSold"] - df["YearBuilt"]
df["RemodAge"]    = df["YrSold"] - df["YearRemodAdd"]
\end{verbatim}

Log-transforms reduzem a skewness de variáveis de área, ajudando
especialmente o meta-modelo linear. As features de idade capturam
depreciação e o efeito de reformas recentes.
\section{Encoding de Variáveis Categóricas}
\label{sec:org9561159}

As 44 variáveis categóricas são transformadas com \textbf{Target Encoding}
dentro de cada fold:

\begin{verbatim}
enc = ce.TargetEncoder(cols=cat_cols, smoothing=10)
enc.fit(X_tr, y_tr)
\end{verbatim}

O parâmetro \texttt{smoothing=10} é crítico: mistura a média da categoria com
a média global proporcionalmente ao tamanho da categoria, evitando
data leakage em categorias raras. O valor 10 (versus 0.3 do baseline)
reduz significativamente o overfitting do encoding.

O encoder é sempre fitado apenas nos dados de treino do fold e
aplicado em validação e teste, respeitando o isolamento de
informação da validação cruzada.
\section{Normalização}
\label{sec:org5974f3e}

As variáveis numéricas são normalizadas com \texttt{RobustScaler} em vez de
\texttt{StandardScaler}:

\begin{verbatim}
scaler = RobustScaler()
\end{verbatim}

\texttt{RobustScaler} usa mediana e IQR em vez de média e desvio padrão,
tornando-o insensível a outliers remanescentes após a filtragem
inicial.
\section{Modelos Base}
\label{sec:orgdfff7f0}

Cinco modelos de gradient boosting com hiperparâmetros distintos
formam a camada base do ensemble. A diversidade entre eles é
intencional: erros não correlacionados se cancelam no stacking.

\begin{center}
\begin{tabular}{lrrrl}
Modelo & \texttt{n\_estimators} & \texttt{learning\_rate} & \texttt{num\_leaves\textasciitilde{}/\textasciitilde{}depth} & Característica\\
\hline
LGB1 & 6000 & 0.005 & 31 & Conservador, balanceado\\
LGB2 & 8000 & 0.003 & 15 & Raso, muito regularizado\\
LGB3 & 5000 & 0.008 & 63 & Muitas folhas, agressivo\\
Cat & 5000 & 0.007 & 6 & Melhor OOF individual\\
XGB & 5000 & 0.005 & 5 & Complementa LGB/Cat\\
\end{tabular}
\end{center}

OOF RMSE individuais (seed 42):
\begin{itemize}
\item LGB1: 0.126874
\item LGB2: 0.123090
\item LGB3: 0.119819
\item CatBoost: \textbf{0.116575} (melhor)
\item XGBoost: 0.125426
\end{itemize}
\section{Validação Cruzada (Out-of-Fold)}
\label{sec:org4bc25e8}

\begin{verbatim}
kf = KFold(n_splits=10, shuffle=True, random_state=42)
\end{verbatim}

10 folds foram utilizados em vez dos 5 do baseline por dois motivos:
cada modelo é avaliado em 10 conjuntos distintos, reduzindo a
variância da estimativa de RMSE; e cada modelo de treino usa 90\% dos
dados em vez de 80\%, gerando predições mais fortes.

As predições Out-of-Fold (OOF) são usadas como features para o
meta-modelo, garantindo que o meta-modelo nunca veja predições feitas
sobre dados em que o modelo base foi treinado.
\section{Meta-Modelo (Stacking)}
\label{sec:org3c95745}

As predições OOF dos cinco modelos base formam uma matriz de entrada
para dois meta-modelos:

\begin{verbatim}
# Meta 1
meta_ridge = RidgeCV(alphas=np.logspace(-4, 2, 50), cv=10)
meta_ridge.fit(X_meta, y)   # RMSE: 0.113588

# Meta 2
meta_en = ElasticNet(alpha=0.0003, l1_ratio=0.5, max_iter=50000)
meta_en.fit(X_meta, y)      # RMSE: 0.114635

# Blend final (50/50)
final = 0.5 * meta_ridge.predict(...) + 0.5 * meta_en.predict(...)
# RMSE: 0.113965
\end{verbatim}

\texttt{RidgeCV} seleciona automaticamente o melhor \texttt{alpha} entre 50 valores
via cross-validation interna. O blend 50/50 com ElasticNet reduz
ligeiramente a variância da predição final em relação a usar apenas um
dos dois.
\section{Resultado Final}
\label{sec:org95c64c8}

\begin{verbatim}
★ OOF RMSE FINAL (blend meta): 0.113965
★ OOF RMSE FINAL (média direta): 0.119427
→ Usando blend de metas (melhor)
\end{verbatim}

A submissão contém os valores preditos diretamente em escala
\texttt{log1p(SalePrice)}, que é o formato esperado pela competição.
\section{Resumo das Decisões de Design}
\label{sec:orgf9a4b72}

\begin{center}
\begin{tabular}{lll}
Decisão & Alternativa descartada & Motivo\\
\hline
Imputação contextual por domínio & Mediana global & NaN em GarageArea significa 0, não mediana\\
RobustScaler & StandardScaler & Mais estável com outliers remanescentes\\
TargetEncoding smoothing=10 & smoothing=0.3 & Menos leakage em categorias raras\\
10-fold CV & 5-fold & Menor variância, mais dados por treino\\
5 modelos heterogêneos & Apenas LightGBM & Erros não correlacionados melhoram o blend\\
RidgeCV como meta & ElasticNet puro & Seleção automática de alpha, mais estável\\
Predição em log scale & \texttt{expm1} antes de submeter & Formato exigido pela competição\\
\end{tabular}
\end{center}
\end{document}
